\chapter{Elementos sobre a economia política da energia}
\section{\textit{Machinery Question}}
Desde meados do século XVIII até o começo do século XIX, estabeleceu-se na Inglaterra a assim chamada \textit{Machinery Question}. Os ares da época são bem explicados por uma passagem de Maxine Berg, citada por \textcite{Pasquinelli2023}, que vale uma citação mais longa:
\begin{quote}
    ``A questão [\textit{Machinery Question}] era central nas relações cotidianas entre patrão e empregado, mas também era de grande interesse teórico e ideológico. A própria tecnologia na base da economia e da sociedade era uma plataforma de desafios e lutas. A máquina era debatida longamente em todos os setores da sociedade. Provocava tanto o clérigo da aldeia quanto o intelectual cosmopolita; preocupava tanto o político quanto o operário e o empregador; o reformador social tanto quanto o cientista e o inventor. Esses grupos disputavam os custos e benefícios da nova tecnologia. Elogiavam a libertação que ela proporcionava das limitações ao crescimento, mas discordavam sobre o impacto que teria sobre os salários, o emprego e a perícia [\textit{wages, employment, and skill}]. Especulavam sobre as mudanças que a máquina traria para as relações sociais, ora acolhendo-as, ora temendo-as. Até mesmo as origens e a propriedade das máquinas foram questionadas. Havia excitação e temor diante dessa força desconhecida que avançava implacavelmente, lançando a velha sociedade fora em seu rastro [\textit{casting the old society in its wake}].'' (Berg \textit{apud} Pasquinelli, 2023, p. 80-1)
\end{quote}

A disrupção social que a maquinaria causou, embora decididamente importante, não cabe à discussão aqui proposta; o que cabe aqui é como a maquinaria afetou os processos de produção capitalista, e como ela não só elucidou, como viabilizou a generalização da forma-mercadoria.

---

No primeiro livro d'O Capital, Marx discorre sobre o que chamou de ``subsunção formal'' e ``subsunção real'' do trabalho pelo capital, que são formas qualitativamente diferentes em que o trabalhador coloca-se perante seu processo de trabalho.

\begin{quote}
    ``Constitui um produto da divisão manufatureira do trabalho opor-lhes [aos trabalhadores] as potências intelectuais do processo material como \textit{propriedade alheia} e como \textit{poder que os domina}. Esse processo de cisão começa na cooperação simples, em que o capitalista representa diante dos trabalhadores individuais a unidade e a vontade do corpo social de trabalho, desenvolve-se na manufatura, que mutila o trabalhador, fazendo dele um trabalhador \textit{parcial}, e se completa na grande indústria, que separa do trabalho a ciência como \textit{potência autônoma de produção} e a obriga a servir ao capital.'' \parencite[p. 435, grifo meu]{Marx2017}
\end{quote}


\# Subsunção real do trabalho, máquina-ferramenta etc


No livro I d'O Capital, Marx separa a caracterização de uma máquina em ``três partes essencialmente distintas'': a máquina-motriz, que ``atua como a força motora do mecanismo inteiro''; o mecanismo de transmissão, que ``regula o movimento, modifica sua forma onde é necessário [...] e o distribui e transmite à máquina-ferramenta''; e, por fim, a máquina-ferramenta, cujo movimento advém de ambas as partes acima, e por meio da qual a máquina ``se apodera do objeto de trabalho e o modifica conforme a uma finalidade'' \parencite[p. 446-7]{Marx2017}.



\begin{quote}
    ``A grande indústria teve, pois, de se apoderar de seu meio característico de produção, a própria máquina, e produzir máquinas por meio de máquinas. Somente assim ela criou sua base técnica adequada e se firmou sobre seus próprios pés. Com a crescente produção mecanizada das primeiras décadas do século XIX, a maquinaria se apoderou gradualmente da fabricação de máquinas-ferramentas. No entanto, foi apenas nas últimas décadas que a colossal construção de ferrovias e a navegação oceânica a vapor deram à luz as ciclópicas máquinas empregadas na construção dos primeiros motores.'' \parencite[p. 458]{Marx2017}
\end{quote}

\# Trabalho como requerendo conhecimento (Thompson) \textcite[p. 85]{Pasquinelli2023}

\# Maquinaria não ``libera trabalho'', e sim \textit{desmembra o processo de trabalho} e, portanto, \textit{cria novos tipos de trabalho ao mesmo tempo em que destroi outros}


- ``machinery did not displace labour. Rather, it differentiated this labour by *dismembering the old craft*.'' (Berg *apud* Pasquinelli, op. cit., p. 89, grifo meu)


\section{A automatização do trabalho e a cisão energia-trabalho}

Não é uma mera coincidência que as medidas de energia na Física moderna têm termos similares aos da produção econômica \parencite{Pasquinelli2022}. \textit{Trabalho}\footnote{Em inglês, \textit{work}; em alemão, \textit{Arbeit}; em francês, \textit{travail}. Os mesmíssimos termos que aparecem na Economia são empregados (em outro contexto, evidentemente) na Física.} trata-se (\textit{grosso modo}) do emprego de uma força $F$ ao longo de alguma distância $d$, processo através do qual transfere-se energia de um meio a outro; as unidades físicas tanto do trabalho quanto da energia são as mesmas, medidas pelo Sistema Internacional (SI) em Joules (J). \textit{Potência}\footnote{Em inglês, \textit{power}; em francês, \textit{puissance} (sinônimo de ``poder''). Em alemão, \textit{Leistung}, que é também empregado no sentido de \textit{eficiência}; inclusive, é o termo que se usa para falar da ``eficiência de uma máquina'' --- não é difícil imaginar esta quantidade física historicamente sendo empregada como um \textit{proxy} da eficiência termodinâmica de motores a vapor no começo do século XIX.} trata-se desta transferência ao longo do tempo, cujas unidades são medidas pelo SI em Watts (W). Estas medidas relacionadas a medidas de energia e de sua mudança tornam-se quintessenciais quando da Revolução Industrial, em que a mensuração mais aderente possível dos meios de produção --- dentre os quais faziam parte, agora, também combustíveis, carvão, óleo etc. --- tornava-se parte crucial da contabilidade do capitalista industrial. 

Não é ao acaso, no mesmo ínterim, que o avanço da maquinaria no Reino Unido tenha causado uma cisão entre o ``trabalho'' humano e o de uma máquina num mesmo processo de produção: o primeiro torna-se \textit{labour}, o segundo torna-se \textit{work}, ambos tendo seu ponto de contato na máquina \textit{in actu} \parencite[p. 11]{Pasquinelli2022}. Como o capitalista ``insiste em obter o que é seu'' e ``[n]ão quer ser furtado'' no que tange a seus dispêndios adiantados de capital\footnote{\textcite[p. 272]{Marx2017}.}, a mensuração mais aderente possível do processo de trabalho, em todas as facetas possíveis, torna-se um \textit{sine qua non} para sua atividade, em particular tornando-se força férrea da concorrência\footnote{A noção de ``concorrência'' em Marx tem um significado mais sutil do que sua acepção usual na Economia. Enquanto esta última tende a abordá-la como ``a quantidade de firmas em um setor de mercado'', Marx utiliza este termo para descrever a ``força social'' que --- não obstante advindo, em última instância, do fazer dos próprios capitais individuais --- acaba por impor-se a eles como ``força da natureza''. [\textbf{Descrição bem rasa; cabe ver se é cabível explicá-lo e, se sim, fazê-lo mais detidamente. Não cabe apenas numa nota de rodapé.}] É neste sentido, por exemplo, que Marx e Engels escreveram, no Manifesto Comunista, que ``Devido ao rápido aperfeiçoamento dos instrumentos de produção e ao constante progresso dos meios de comunicação, a burguesia arrasta para a torrente da civilização mesmo as nações mais bárbaras. Os baixos preços de seus produtos são a artilharia pesada que destrói todas as muralhas da China e obriga a capitularem os bárbaros mais tenazmente hostis aos estrangeiros.''.}. Dessa forma, Pasquinelli observa que

\begin{quote}
    ``as máquinas são inventadas pela organização do trabalho e, posteriormente, elas [as máquinas] acabam impondo ao trabalho suas próprias métricas. Métricas de trabalho e automação do trabalho parecem estar intimamente ligados.'' \parencite[p. 7]{Pasquinelli2022}
\end{quote}

